% Master's thesis -- UMONS, Faculté des Sciences
% Based on the UMONS master's thesis template (A. Buys), converted to UTF-8 and English.
\documentclass[a4paper,11pt,notitlepage]{report}

\usepackage[utf8]{inputenc}
\usepackage[T1]{fontenc}
\usepackage{mathptmx} % Times, as in the UMONS template
\usepackage[english]{babel}
\usepackage{amsmath,amssymb,amsthm}
\usepackage{graphicx}
\usepackage[vmargin=2.6cm,hmargin=2.9cm]{geometry}
\usepackage{indentfirst}
\usepackage{xcolor}
\usepackage{listings}
\usepackage{url}
\usepackage[hidelinks]{hyperref}
\usepackage[capitalise,noabbrev]{cleveref}
\usepackage{booktabs}
\usepackage{tabularx}

% ---------------------------------------------------------------------------
% YAML listings
% ---------------------------------------------------------------------------
\lstdefinelanguage{yaml}{
  keywords={true,false,null},
  comment=[l]{\#},
  morestring=[b]', morestring=[b]",
  sensitive=true
}
\lstset{
  language=yaml,
  basicstyle=\ttfamily\scriptsize,
  columns=fullflexible,
  keepspaces=true,
  numbers=left, numberstyle=\tiny\color{gray},
  frame=single, rulecolor=\color{gray!50},
  commentstyle=\color{gray}\itshape,
  stringstyle=\color{teal!70!black},
  breaklines=true,
  showstringspaces=false,
  captionpos=b
}
% Path to the exemplar workflows (shared with the test suite)
\newcommand{\ex}[1]{../examples/#1}

% ---------------------------------------------------------------------------
% Title page
% ---------------------------------------------------------------------------
\pagestyle{myheadings}

\newcommand{\STitre}{UNIVERSIT\'E DE MONS}
\newcommand{\FTitre}{FACULT\'E DES SCIENCES}
\newcommand{\Titre}{A Refactoring Tool for GitHub Actions Workflows}

\theoremstyle{definition}
\newtheorem{defn}{Definition}
\newtheorem{exa}{Example}

\title{
  \begin{minipage}\linewidth\centering
  \LARGE\bf\STitre
  \vskip20pt
  \Large\bf\FTitre
  \vskip20pt
  \Huge \Titre
  \end{minipage}
}
\author{Abdellah BOUMAIZA\\[1.5ex]
Master en Sciences Informatiques\\[1.5ex]
{\tt \small abdellah.boumaiza@student.umons.ac.be} % TODO check e-mail
}
\date{Academic year 2026--2027}

\begin{document}

\maketitle
\markright{\Titre}
\thispagestyle{empty}

\vfill
\begin{quote}
Master's thesis carried out in the \textbf{Software Engineering Lab}, % TODO exact service name
as part of the \textbf{Master en Sciences Informatiques}.
\end{quote}
\begin{quote}
\textbf{Supervisor:} Prof.\ Tom MENS \\
\textbf{Co-supervisor:} Aref TALEBZADEH BARDSIRI
\end{quote}
\vfill

\includegraphics[width=0.2\textwidth]{img/FS.pdf}\hfill\includegraphics[width=0.2\textwidth]{img/logo_UMONS.pdf}

\newpage
The author, Abdellah BOUMAIZA, declares having complied with the ethical rules in force,
including the University's charter on the use of Artificial Intelligence.

\section*{Use of AI tools}
% TODO: describe honestly which tools were used and for what. Example:
Generative AI tools (Claude, Anthropic) were used to help structure the document,
draft the \LaTeX{} skeleton and continuous-integration workflow of this repository,
propose a first version of the refactoring exemplars, and rephrase parts of the text.
All content was reviewed, checked against the cited sources and is the responsibility of the author.

\newpage
\tableofcontents
\newpage

\setlength{\parindent}{12pt}
\setlength{\parskip}{1ex plus 0.5ex minus 0.2ex}

% ===========================================================================
\chapter*{Introduction}\addcontentsline{toc}{chapter}{Introduction}
\label{ch:intro}
% ===========================================================================

% TODO: paste/adapt the first text sent to the supervisor here.

GitHub Actions has become the dominant continuous-integration and automation service
for repositories hosted on GitHub~\cite{Decan2022,Kinsman2021}.
Workflows are YAML files that tend to grow and to be duplicated within and across
repositories. As for source code, this duplication hampers maintenance and can be
addressed by \emph{refactoring}, i.e., behaviour-preserving transformations that improve
the structure of an artefact~\cite{Fowler2018,MensTourwe2004}.

This thesis focuses on refactorings that improve \emph{reuse} in workflows.
Its goals are (i) to catalogue the reuse mechanisms offered by GitHub Actions and the
refactorings they enable (\cref{ch:mechanisms,ch:catalogue}),
and (ii) to build a tool that detects refactoring opportunities and applies them automatically.

% ===========================================================================
\chapter{Reuse mechanisms in GitHub Actions}
\label{ch:mechanisms}
% ===========================================================================

GitHub Actions offers several mechanisms to avoid duplicating workflow
configuration~\cite{GHDocsReuse}. They differ in the \emph{granularity} of what is reused
(a YAML node, a sequence of steps, a job, a whole workflow), in their \emph{scope}
(one file, one repository, an organisation, the whole platform) and in the moment at which
reuse happens (parse time, run time, or creation time).
\Cref{tab:mechanisms} summarises them.

\begin{table}[ht]
\centering\small
\begin{tabularx}{\textwidth}{@{}lXXl@{}}
\toprule
\textbf{Mechanism} & \textbf{Reused unit} & \textbf{Scope} & \textbf{When}\\
\midrule
YAML anchors/aliases~\cite{GHChangelogAnchors2025,YAML122} & any YAML node & one file & parse time\\
Composite action~\cite{GHDocsComposite} & sequence of steps & repo / public & run time\\
JavaScript / Docker action & one step (custom code) & repo / public & run time\\
Reusable workflow & one or more jobs & repo / org / public & run time\\
Workflow template~\cite{GHChangelogAnchors2025} & whole workflow file & organisation & creation time\\
Matrix strategy & one job, parametrised & one job & run time\\
Chained workflows (\texttt{workflow\_run}) & a set of jobs & one repository & run time\\
\bottomrule
\end{tabularx}
\caption{Reuse mechanisms available in GitHub Actions.}
\label{tab:mechanisms}
\end{table}

% TODO: one section per mechanism (syntax, constraints, limits).

% ===========================================================================
\chapter{A catalogue of reuse refactorings}
\label{ch:catalogue}
% ===========================================================================

Each refactoring is illustrated by a pair of artificial but realistic workflows
(\emph{before} / \emph{after}) stored in the \texttt{examples/} folder of the repository.
These exemplars form a small ground truth: they document the catalogue and will serve as a
test suite for the refactoring tool. All of them are validated by
\texttt{actionlint}, and, where the refactoring is a pure syntactic rewriting,
an automatic equivalence check is provided (\texttt{examples/check.py}).

\begin{table}[ht]
\centering\small
\begin{tabularx}{\textwidth}{@{}llX@{}}
\toprule
\textbf{Id} & \textbf{Refactoring} & \textbf{Strictly behaviour-preserving?}\\
\midrule
R1 & Extract Composite Action & yes\\
R2 & Extract Reusable Workflow & yes (modulo UI/log structure)\\
R3 & Introduce YAML Anchors & yes (identical after alias expansion)\\
R4 & Extract Workflow Template & yes for the instantiated file; no propagation\\
R5 & Split Workflow (\texttt{workflow\_run}) & no: default-branch version, security context\\
R6 & Introduce Matrix & yes with \texttt{fail-fast: false}; job ids change\\
\bottomrule
\end{tabularx}
\caption{Refactorings of the catalogue.}
\label{tab:catalogue}
\end{table}

% ---------------------------------------------------------------------------
\section{R1 -- Extract Composite Action}\label{sec:r1}

\textbf{Smell.} The same sequence of steps is repeated in several jobs.
\textbf{Mechanism.} Composite action~\cite{GHDocsComposite}.

\lstinputlisting[caption={R1, before: three jobs repeat the Node setup.},label=lst:r1-before]{\ex{R1-extract-composite-action/before/.github/workflows/ci.yml}}
\lstinputlisting[caption={R1, after: extracted composite action.}]{\ex{R1-extract-composite-action/after/.github/actions/setup-node-project/action.yml}}
\lstinputlisting[caption={R1, after: the workflow calls the local action.}]{\ex{R1-extract-composite-action/after/.github/workflows/ci.yml}}

\textbf{Pitfalls.} \texttt{run} steps must declare a \texttt{shell}; a local action is only
available after \texttt{actions/checkout}; job-level keys (\texttt{runs-on},
\texttt{services}, \texttt{permissions}) cannot be extracted.

% ---------------------------------------------------------------------------
\section{R2 -- Extract Reusable Workflow}\label{sec:r2}

\textbf{Smell.} The same \emph{job} is duplicated in several workflows.
\textbf{Mechanism.} Reusable workflow (\texttt{on: workflow\_call})~\cite{GHDocsReuse}.

\lstinputlisting[caption={R2, before: staging deployment (production is identical up to 5 values).}]{\ex{R2-extract-reusable-workflow/before/.github/workflows/deploy-staging.yml}}
\lstinputlisting[caption={R2, after: the reusable workflow.}]{\ex{R2-extract-reusable-workflow/after/.github/workflows/_deploy.yml}}
\lstinputlisting[caption={R2, after: a caller.}]{\ex{R2-extract-reusable-workflow/after/.github/workflows/deploy-staging.yml}}

% ---------------------------------------------------------------------------
\section{R3 -- Introduce YAML Anchors}\label{sec:r3}

\textbf{Smell.} Repeated YAML nodes \emph{within one file}.
\textbf{Mechanism.} YAML anchors and aliases~\cite{YAML122}, supported by GitHub Actions
since September 2025~\cite{GHChangelogAnchors2025}.

\lstinputlisting[caption={R3, after (the \emph{before} version, with the three blocks inlined three times, is in the repository).}]{\ex{R3-introduce-yaml-anchors/after/.github/workflows/python-ci.yml}}

\textbf{Remark.} Anchors are expanded at parse time, so behaviour preservation can be checked
mechanically by comparing the parsed documents. Anchors are file-local and not parametrisable.

% ---------------------------------------------------------------------------
\section{R4 -- Extract Workflow Template}\label{sec:r4}

\textbf{Smell.} The same workflow file is copy-pasted across repositories of an organisation.
\textbf{Mechanism.} Organisation workflow template, possibly in a non-public
\texttt{.github} repository since 2025~\cite{GHChangelogAnchors2025}.

\lstinputlisting[caption={R4, the template (\texttt{\$default-branch} is a placeholder).}]{\ex{R4-extract-workflow-template/after/org-dotgithub/workflow-templates/java-maven-ci.yml}}
\lstinputlisting[caption={R4, template metadata.},language={}]{\ex{R4-extract-workflow-template/after/org-dotgithub/workflow-templates/java-maven-ci.properties.json}}

\textbf{Remark.} Unlike R1--R3, a template is \emph{copied} at creation time: duplication is
standardised rather than removed, and later changes are not propagated.
Combining R4 with R2 (a template that only calls a reusable workflow) gives both.

% ---------------------------------------------------------------------------
\section{R5 -- Split Workflow}\label{sec:r5}

\textbf{Smell.} A long sequential workflow mixes several concerns.
\textbf{Mechanism.} Two regular workflows chained with the \texttt{workflow\_run} event.

\lstinputlisting[caption={R5, before.}]{\ex{R5-split-workflow/before/.github/workflows/ci-and-docs.yml}}
\lstinputlisting[caption={R5, after: the second workflow.}]{\ex{R5-split-workflow/after/.github/workflows/docs.yml}}

\textbf{Remark.} This refactoring is \emph{not} strictly behaviour-preserving: the triggered
workflow always uses the default-branch version of its file and runs in a privileged context,
even for runs triggered by pull requests from forks.

% ---------------------------------------------------------------------------
\section{R6 -- Introduce Matrix}\label{sec:r6}

\textbf{Smell.} Jobs that differ only by a few values.
\textbf{Mechanism.} \texttt{strategy.matrix}.

\lstinputlisting[caption={R6, after (the \emph{before} version contains nine near-identical jobs).}]{\ex{R6-introduce-matrix/after/.github/workflows/tests.yml}}

% ===========================================================================
\chapter{Conclusion}\label{ch:conclusion}
% ===========================================================================

% TODO

% ===========================================================================
\bibliographystyle{plain}
\bibliography{references}
\addcontentsline{toc}{chapter}{Bibliography}

\end{document}
